\documentclass[12pt]{article}
\usepackage[utf8]{inputenc}
\usepackage{graphicx}
\usepackage{subcaption}
\usepackage{amsmath}
\usepackage{fancyhdr}
\usepackage{geometry}
\usepackage{float}
\usepackage{listings}
\usepackage{xcolor}
\usepackage{array}
\usepackage{caption}
\usepackage{booktabs}
\usepackage[english]{babel}
\usepackage{amssymb}
\usepackage{hyperref}

\lstset{
    basicstyle=\ttfamily\footnotesize,
    backgroundcolor=\color{gray!10},
    keywordstyle=\color{blue},
    commentstyle=\color{green!50!black},
    stringstyle=\color{red},
    numbers=left,
    numberstyle=\tiny\color{gray},
    stepnumber=1,
    numbersep=10pt,
    showstringspaces=false,
    frame=single,
    tabsize=4,
    breaklines=true,
    breakatwhitespace=false,
    language=Python
}

\linespread{1.25}
\setlength{\parindent}{0.8cm}
\setlength{\parskip}{0em}
\renewcommand{\headrulewidth}{0pt}
\geometry{a4paper, portrait, margin=1in}
\setlength{\headheight}{14.49998pt}
\graphicspath{{Lab7/charts/}}

\begin{document}

% ─── TITLE PAGE ───────────────────────────────────────────────────────────────
\begin{titlepage}
  \begin{center}
    \textbf{\large Ministry of Education of Republic of Moldova}\\[0.1cm]
    \textbf{\large Technical University of Moldova}\\[0.1cm]
    \textbf{\large Faculty of Computers, Informatics and Microelectronics}\\[1.2cm]

    \vspace{40mm}

    \textsc{\Large Laboratory Work 7:}\\[0.2cm]
    \textsc{\Large Empirical Analysis of Greedy Algorithms}\\[0.2cm]
    \textsc{\large Prim and Kruskal Minimum Spanning Tree Algorithms}\\[0.4cm]

    \vspace{65mm}

    \begin{minipage}[t]{0.4\textwidth}
      \begin{flushleft}\large
        \emph{Author:}\\
        st. gr. FAF-242\\[0.4cm]
        \emph{Verified:}\\
        asist. univ.
      \end{flushleft}
    \end{minipage}
    ~
    \begin{minipage}[t]{0.4\textwidth}
      \begin{flushright}\large
        \emph{}\\
        Tacu Igor\\[0.4cm]
        \emph{}\\
        Fi\c{s}tic Cristofor
      \end{flushright}
    \end{minipage}\\[2cm]

    \vspace{5mm}
    \large Chi\c{s}in\u{a}u -- 2026

    \vfill
  \end{center}
\end{titlepage}

% ─── TABLE OF CONTENTS ────────────────────────────────────────────────────────
\clearpage
\thispagestyle{empty}
\tableofcontents

\clearpage
\setcounter{page}{2}
\pagestyle{fancy}
\fancyhf{}
\rhead{\thepage}

% ─── 1. OBJECTIVE ─────────────────────────────────────────────────────────────
\section*{OBJECTIVE}
\addcontentsline{toc}{section}{OBJECTIVE}

\hspace{0.8cm} The objective of this laboratory work is to study the greedy algorithm
design technique and apply it to the Minimum Spanning Tree (MST) problem.
Specifically, the goals are:

\begin{enumerate}
  \item Study the greedy algorithm design technique.
  \item Implement Prim's and Kruskal's MST algorithms in Python.
  \item Conduct empirical analysis of both algorithms.
  \item Investigate how increasing the number of nodes affects execution time
        across sparse, medium, and dense graph types.
  \item Present the results graphically and draw conclusions.
\end{enumerate}

% ─── 2. THEORETICAL BACKGROUND ────────────────────────────────────────────────
\section*{THEORETICAL BACKGROUND}
\addcontentsline{toc}{section}{THEORETICAL BACKGROUND}

\subsection*{Greedy Algorithms}
\addcontentsline{toc}{subsection}{Greedy Algorithms}

\hspace{0.8cm} A \textbf{greedy algorithm} builds a solution incrementally by always
making the locally optimal choice at each step, hoping that these local choices will
lead to a globally optimal solution. Greedy algorithms are characterised by:

\begin{itemize}
  \item \textbf{Greedy choice property:} a globally optimal solution can be assembled
        from locally optimal choices made independently.
  \item \textbf{Optimal substructure:} an optimal solution to the problem contains
        optimal solutions to its sub-problems.
  \item \textbf{No backtracking:} once a choice is made it is never reconsidered.
\end{itemize}

\hspace{0.8cm} Classical examples include activity selection, fractional knapsack,
Huffman coding, and MST algorithms. Not every problem admits a greedy solution;
correctness must be proved via exchange arguments or matroid theory.

\subsection*{Minimum Spanning Tree}
\addcontentsline{toc}{subsection}{Minimum Spanning Tree}

\hspace{0.8cm} Given a connected, undirected, weighted graph $G = (V, E, w)$, a
\textbf{spanning tree} is an acyclic subgraph that connects every vertex.
A \textbf{Minimum Spanning Tree (MST)} is a spanning tree whose total edge weight

\[
  W(T) = \sum_{(u,v) \in T} w(u,v)
\]

is minimised. Every connected graph has at least one MST; it is unique when all
edge weights are distinct. The MST has exactly $|V| - 1$ edges.

\subsection*{Prim's Algorithm}
\addcontentsline{toc}{subsection}{Prim's Algorithm}

\hspace{0.8cm} Prim's algorithm grows the MST one vertex at a time, starting from
an arbitrary source. At each step the cheapest edge crossing the cut between the
current MST and the remaining vertices is added. Using a binary min-heap (priority
queue) the algorithm achieves $O((V + E)\log V)$ time.

\begin{enumerate}
  \item Initialise a key array: $\text{key}[s] = 0$, $\text{key}[v] = \infty$ for all $v \neq s$.
  \item Insert $(0, s)$ into the min-heap.
  \item While the heap is non-empty:
  \begin{enumerate}
    \item Extract the vertex $u$ with the minimum key.
    \item Mark $u$ as in the MST; record the edge from its parent.
    \item For each neighbour $v$ of $u$ not yet in the MST: if $w(u,v) < \text{key}[v]$,
          update $\text{key}[v]$ and push $(w(u,v), v)$ onto the heap.
  \end{enumerate}
\end{enumerate}

\textbf{Complexity:}
\begin{itemize}
  \item Time: $O((V + E)\log V)$ with a binary heap; $O(E + V\log V)$ with a Fibonacci heap.
  \item Space: $O(V)$ for the key array and heap.
\end{itemize}

\subsection*{Kruskal's Algorithm}
\addcontentsline{toc}{subsection}{Kruskal's Algorithm}

\hspace{0.8cm} Kruskal's algorithm takes a global view: it sorts all edges by weight
and greedily adds each edge to the MST provided it does not create a cycle. The
cycle check is performed in near-constant time using a \textbf{Union-Find} (disjoint
set union) data structure with path compression and union by rank.

\begin{enumerate}
  \item Sort all edges $E$ in non-decreasing order of weight: $O(E \log E)$.
  \item Initialise a Union-Find structure: each vertex is its own component.
  \item For each edge $(u, v, w)$ in sorted order:
  \begin{enumerate}
    \item If $\text{find}(u) \neq \text{find}(v)$: add $(u, v, w)$ to the MST and call $\text{union}(u, v)$.
    \item Otherwise: skip (adding would create a cycle).
  \end{enumerate}
  \item Stop when $|V| - 1$ edges have been added.
\end{enumerate}

\textbf{Complexity:}
\begin{itemize}
  \item Time: $O(E \log E)$ dominated by the sort; Union-Find operations cost
        $O(E \cdot \alpha(V)) \approx O(E)$ where $\alpha$ is the inverse Ackermann function.
  \item Space: $O(E)$ for the sorted edge list plus $O(V)$ for the Union-Find structure.
\end{itemize}

\subsection*{Comparison of Algorithms}
\addcontentsline{toc}{subsection}{Comparison of Algorithms}

\begin{center}
\begin{tabular}{lll}
  \toprule
  \textbf{Property}      & \textbf{Prim}              & \textbf{Kruskal}         \\
  \midrule
  Strategy               & Vertex-centric (grow MST)  & Edge-centric (sort \& add) \\
  Time complexity        & $O((V+E)\log V)$           & $O(E \log E)$            \\
  Space complexity       & $O(V)$                     & $O(E)$                   \\
  Best for               & Dense graphs               & Sparse graphs            \\
  Handles negative $w$   & Yes                        & Yes                      \\
  Data structure         & Min-heap                   & Union-Find               \\
  \bottomrule
\end{tabular}
\end{center}

% ─── 3. IMPLEMENTATION ────────────────────────────────────────────────────────
\section*{IMPLEMENTATION}
\addcontentsline{toc}{section}{IMPLEMENTATION}

\hspace{0.8cm} Both algorithms are implemented in Python using \texttt{networkx}
undirected graphs. The full source is in \texttt{Lab7/src/}.

\subsection*{Prim's Algorithm}
\addcontentsline{toc}{subsection}{Prim's Algorithm}

\begin{lstlisting}[caption={Prim's algorithm using a binary min-heap}, label={lst:prim}]
def prim(graph, source: int = 0) -> dict:
    in_mst = set()
    key = {node: math.inf for node in graph.nodes()}
    key[source] = 0
    heap = [(0, source, -1)]   # (weight, node, parent)
    mst_edges = []
    total_weight = 0.0
    max_heap = 1

    t0 = time.perf_counter()
    while heap:
        w, u, parent = heapq.heappop(heap)
        if u in in_mst:
            continue
        in_mst.add(u)
        if parent != -1:
            mst_edges.append((parent, u, w))
            total_weight += w
        for v, data in graph[u].items():
            if v not in in_mst:
                edge_w = data.get("weight", 1)
                if edge_w < key[v]:
                    key[v] = edge_w
                    heapq.heappush(heap, (edge_w, v, u))
                    if len(heap) > max_heap:
                        max_heap = len(heap)
    elapsed_ms = (time.perf_counter() - t0) * 1_000
    return {"time_ms": elapsed_ms, "mst_edges": mst_edges,
            "total_weight": total_weight, "memory_peak": max_heap}
\end{lstlisting}

\subsection*{Kruskal's Algorithm}
\addcontentsline{toc}{subsection}{Kruskal's Algorithm}

\begin{lstlisting}[caption={Kruskal's algorithm with Union-Find}, label={lst:kruskal}]
def kruskal(graph) -> dict:
    edges = sorted([(data.get("weight", 1), u, v)
                    for u, v, data in graph.edges(data=True)])
    parent = {node: node for node in graph.nodes()}
    rank   = {node: 0    for node in graph.nodes()}

    def find(x):
        while parent[x] != x:
            parent[x] = parent[parent[x]]  # path halving
            x = parent[x]
        return x

    def union(x, y):
        px, py = find(x), find(y)
        if px == py:
            return False
        if rank[px] < rank[py]:
            px, py = py, px
        parent[py] = px
        if rank[px] == rank[py]:
            rank[px] += 1
        return True

    mst_edges = []
    total_weight = 0.0
    t0 = time.perf_counter()
    for w, u, v in edges:
        if union(u, v):
            mst_edges.append((u, v, w))
            total_weight += w
    elapsed_ms = (time.perf_counter() - t0) * 1_000
    return {"time_ms": elapsed_ms, "mst_edges": mst_edges,
            "total_weight": total_weight, "memory_peak": len(edges)}
\end{lstlisting}

% ─── 4. ALGORITHM ANIMATIONS ──────────────────────────────────────────────────
\section*{ALGORITHM ANIMATIONS}
\addcontentsline{toc}{section}{ALGORITHM ANIMATIONS}

\subsection*{Prim's Algorithm Animation}
\addcontentsline{toc}{subsection}{Prim's Algorithm Animation}

\hspace{0.8cm} The animation runs on a 12-node undirected sparse graph.
Each frame shows the node just added to the MST (red), nodes already in the
MST (green), unvisited nodes (grey), and current MST edges (orange).
Edge weights are displayed on the graph throughout.
Full animation: \texttt{Lab7/charts/prim.gif}.

\begin{figure}[H]
  \centering
  \begin{subfigure}[t]{0.48\textwidth}
    \includegraphics[width=\textwidth]{prim_step00.png}
    \caption{Initial state: source node 0 added to MST}
  \end{subfigure}
  \hfill
  \begin{subfigure}[t]{0.48\textwidth}
    \includegraphics[width=\textwidth]{prim_step04.png}
    \caption{Mid-traversal: several nodes incorporated}
  \end{subfigure}
  \caption{Prim's algorithm on the 12-node demo graph.}
  \label{fig:prim-steps}
\end{figure}

\subsection*{Kruskal's Algorithm Animation}
\addcontentsline{toc}{subsection}{Kruskal's Algorithm Animation}

\hspace{0.8cm} The animation uses the same 12-node graph.
Each frame highlights the edge currently under consideration: purple if accepted
into the MST, red if rejected due to forming a cycle. Accepted MST edges remain
orange. The animation makes the greedy edge-sorting strategy clearly visible.
Full animation: \texttt{Lab7/charts/kruskal.gif}.

\begin{figure}[H]
  \centering
  \begin{subfigure}[t]{0.48\textwidth}
    \includegraphics[width=\textwidth]{kruskal_step00.png}
    \caption{First edge considered (lowest weight)}
  \end{subfigure}
  \hfill
  \begin{subfigure}[t]{0.48\textwidth}
    \includegraphics[width=\textwidth]{kruskal_step05.png}
    \caption{Mid-run: some edges accepted, some rejected}
  \end{subfigure}
  \caption{Kruskal's algorithm on the 12-node demo graph.}
  \label{fig:kruskal-steps}
\end{figure}

% ─── 5. RESULTS AND ANALYSIS ──────────────────────────────────────────────────
\section*{RESULTS AND ANALYSIS}
\addcontentsline{toc}{section}{RESULTS AND ANALYSIS}

\subsection*{Sparse Graphs}
\addcontentsline{toc}{subsection}{Sparse Graphs}

\begin{figure}[H]
  \centering
  \includegraphics[width=\textwidth]{time_sparse.png}
  \caption{Prim vs Kruskal on sparse graphs (uniform weights).}
  \label{fig:time-sparse}
\end{figure}

\hspace{0.8cm} On sparse graphs (Figure~\ref{fig:time-sparse}), both algorithms
are extremely fast even at $n = 2000$. With $E \approx 3V$, the edge sort in Kruskal
costs $O(V \log V)$ and Prim's heap also processes $O(V)$ entries, making both
essentially linear in practice. Kruskal has a slight advantage due to its simpler
inner loop, while Prim incurs overhead from repeated heap insertions.

\subsection*{Dense Graphs}
\addcontentsline{toc}{subsection}{Dense Graphs}

\begin{figure}[H]
  \centering
  \includegraphics[width=\textwidth]{time_dense.png}
  \caption{Prim vs Kruskal on dense graphs (uniform weights).}
  \label{fig:time-dense}
\end{figure}

\hspace{0.8cm} On dense graphs (Figure~\ref{fig:time-dense}), the difference
becomes more pronounced. With $E \approx 0.35 V^2$ (undirected), Kruskal must sort
$O(V^2)$ edges, which dominates its runtime. Prim only ever inserts each vertex once
into the MST and explores at most $E$ edges total, giving it a lower constant factor
when edge processing is cheap. At $n = 500$, both algorithms slow noticeably but
Prim maintains a slight edge.

\subsection*{Weight Distribution Effects}
\addcontentsline{toc}{subsection}{Weight Distribution Effects}

\begin{figure}[H]
  \centering
  \includegraphics[width=\textwidth]{time_by_weights.png}
  \caption{Effect of edge-weight distribution on sparse graph performance.}
  \label{fig:time-weights}
\end{figure}

\hspace{0.8cm} Figure~\ref{fig:time-weights} shows that the weight distribution
(uniform, skewed, or negative) has no meaningful effect on execution time for either
algorithm. Kruskal performs the same number of comparisons during sorting regardless
of value range, and Prim's heap operations are determined by relative order, not
magnitude. Both algorithms handle negative weights correctly because MST correctness
depends only on comparison, not absolute values.

\subsection*{Memory Comparison}
\addcontentsline{toc}{subsection}{Memory Comparison}

\begin{figure}[H]
  \centering
  \includegraphics[width=\textwidth]{memory_comparison.png}
  \caption{Peak memory: Prim's heap size vs Kruskal's sorted edge list.}
  \label{fig:memory}
\end{figure}

\hspace{0.8cm} Figure~\ref{fig:memory} highlights a key practical difference.
Prim's heap holds at most $O(V)$ entries at any point. Kruskal's sorted edge list
holds all $|E|$ edges simultaneously before processing begins. On sparse graphs the
two are comparable ($|E| \approx 3V$). On dense graphs, Kruskal requires $O(V^2)$
memory for the edge list while Prim stays at $O(V)$ — a significant advantage as
$n$ grows.

\subsection*{Full Comparison Grid}
\addcontentsline{toc}{subsection}{Full Comparison Grid}

\begin{figure}[H]
  \centering
  \includegraphics[width=\textwidth]{comparison_all.png}
  \caption{Prim (top row) vs Kruskal (bottom row) across sparse, medium, and dense graphs.}
  \label{fig:comparison-all}
\end{figure}

\hspace{0.8cm} Figure~\ref{fig:comparison-all} provides the complete picture.
On sparse and medium graphs the two algorithms are essentially equivalent in speed.
On dense graphs, Kruskal's $O(E \log E)$ sort becomes the bottleneck as $E$ grows
quadratically, while Prim scales slightly better due to its $O(V)$ memory footprint
and fewer priority queue operations relative to $|E|$.

% ─── 6. CONCLUSION ────────────────────────────────────────────────────────────
\section*{CONCLUSION}
\addcontentsline{toc}{section}{CONCLUSION}

\hspace{0.8cm} Through empirical analysis, Prim and Kruskal were measured across
three graph densities, three weight distributions, and a range of node counts.
Key findings:

\begin{itemize}
  \item \textbf{Both algorithms produce correct MSTs} of equal total weight on all
        tested graphs, including those with negative edge weights.

  \item \textbf{On sparse graphs} both algorithms are near-linear and effectively
        equivalent. Kruskal is marginally faster due to simpler loop structure.

  \item \textbf{On dense graphs} Kruskal's upfront $O(E \log E)$ edge sort becomes
        the bottleneck. Prim's heap-based approach, which processes vertices
        sequentially, scales slightly better.

  \item \textbf{Memory:} Prim uses $O(V)$ space (heap) at all densities. Kruskal
        requires $O(E)$ for the sorted edge list, which grows to $O(V^2)$ on dense
        graphs — a significant disadvantage for large, dense inputs.

  \item \textbf{Weight distribution} has no effect on the execution time of either
        algorithm, as both rely only on relative edge ordering, not absolute values.
\end{itemize}

\textbf{Algorithm selection guidance:}
\begin{itemize}
  \item Use \textbf{Kruskal} for sparse graphs or when an edge list is already
        available: sorting is cheap and the Union-Find overhead is negligible.
  \item Use \textbf{Prim} for dense graphs or when memory is constrained:
        its $O(V)$ space usage is far more efficient than Kruskal's $O(E) = O(V^2)$
        on dense inputs. With a Fibonacci heap, Prim achieves the optimal
        $O(E + V \log V)$ for the MST problem.
\end{itemize}

\vspace{0.5cm}
\noindent\textbf{Repository:} \url{https://github.com/igortacu/AA}

\pagebreak
\end{document}
